% Retired quiz items -- NOT part of any document.
%
% This file is not \input by pre_form_*.tex, post_form_*.tex or key_form_*.tex,
% and it is not compiled. It is an archive, kept so that the items below can be
% reinstated in a later cohort without rewriting them. To reinstate a cell, move
% the A item into part_c_form_a.tex and the B item into part_c_form_b.tex at the
% same position, inside the \begin{questions} environment, and renumber both
% files plus docs/blueprint.md and docs/coding_sheet.csv.
%
% ---------------------------------------------------------------------------
% Why these items were retired
% ---------------------------------------------------------------------------
% On 2026-08-05 the quiz was cut from 28 to 20 items per form, because 28 items
% did not fit the time available in the session. The eight cells below were
% selected so that the surviving 20 keep the properties the instrument relies on:
%
%   - all 13 topics remain represented (T3 and T4, which had three items each,
%     lose exactly one; six topics drop from two items to one),
%   - key positions stay balanced at 5 items each at (a), (b), (c), (d) in both
%     forms -- the retired set contains exactly two of each key letter per form,
%   - the level and difficulty mix stays proportional: the retired set is
%     2 remember / 3 understand / 3 apply-analyse and 2 easy / 3 medium / 3 hard,
%   - three of the four items that docs/blueprint.md listed as needing a second
%     opinion are among the retired ones (see the notes on A21, A24 and A28/B28
%     below).
%
% Item-specific notes carried over from the "Known judgement calls" section of
% docs/blueprint.md, which is why these three were preferred for retirement:
%
%   A21  The most memorisation-flavoured item in the set. Its ordering of
%        isolation levels follows the lab handout; if the TAs teach a different
%        ordering, it has to be rewritten rather than marked wrong.
%   A24  Option (c) -- "deciding whether the request has been fulfilled" -- is
%        not delegable, since it is the coordinating agent's own job, but a
%        student who reads it as "a sub-agent could double-check this" has a
%        defensible position.
%   A28 / B28  Accept several enforcement mechanisms as a pass. Deliberate --
%        the property under test is that the sentinel is never read -- but it
%        makes the items harder than their distractors suggest.
%
% ---------------------------------------------------------------------------
% Renumbering applied to the surviving items
% ---------------------------------------------------------------------------
% The 20 surviving items were renumbered 1-20 and their blueprint cell IDs were
% reassigned to A01-A20 / B01-B20. The retired items below keep their ORIGINAL
% cell IDs (A02/B02 ... A28/B28), which refer to the 28-item version of the
% instrument. The mapping is:
%
%   old  01  02  03  04  05  06  07  08  09  10  11  12  13  14
%   new  01   -  02  03   -  04  05  06   -  07  08  09  10   -
%
%   old  15  16  17  18  19  20  21  22  23  24  25  26  27  28
%   new  11   -  12  13  14  15   -  16  17   -  18  19  20   -
%
% So a reference to A17 in the 28-item version is item A12 in the 20-item
% version; a reference to A16 has no counterpart, because that cell is archived
% here.


% =====================================================================
% Cell 02 -- T1 / understand / medium
% =====================================================================

% --- Form A, item A02 -------------------------------------------------
\qitem A team's harness sends the whole conversation history with every LLM call. To save tokens, a
student proposes sending only the newest user message instead. Why would the agent break?
\begin{choices}
  \choice Because the API rejects requests that fall below a minimum number of tokens.
  \choice Because the earlier calls have already adjusted the model, and the history keeps it consistent.
  \choice Because the sampling temperature becomes unstable when the prompt length varies between calls.
  \CorrectChoice Because the model would no longer see the earlier tool results or its own reasoning, so it could not continue the task.
\end{choices}
\dontknow
\itemmeta{T1 --- LLM role \& statelessness}{understand}{medium}{A02}
\rationale{The history \emph{is} the agent's working state. (b) is the ``the model learns during the run''
misconception.}

% --- Form B, item B02 -------------------------------------------------
\qitem Two calls are made to the same LLM endpoint five minutes apart, within one agent run. If the
harness sends nothing but the new user message in the second call, what does the model still have from the
first one?
\begin{choices}
  \CorrectChoice Nothing at all --- each call is independent, so whatever is not resent is gone.
  \choice The tool schemas, because they were registered with the endpoint on first use.
  \choice The system prompt, because it is cached on the server for the duration of the session.
  \choice The last assistant message, because it is required to continue a conversation.
\end{choices}
\dontknow
\itemmeta{T1 --- LLM role \& statelessness}{understand}{medium}{B02}
\rationale{Same cell as A02: the harness, not the endpoint, owns the state.}


% =====================================================================
% Cell 05 -- T3 / remember / easy
% =====================================================================

% --- Form A, item A05 -------------------------------------------------
\qitem The agentic loop is often described as a state machine. Between which two states does it
alternate?
\begin{choices}
  \choice A training state, in which the model adapts, and an inference state, in which it answers.
  \choice A planning state, in which the task is decomposed, and a reporting state, in which the user is informed.
  \CorrectChoice A reasoning state, in which the LLM is called, and an acting state, in which a tool is called.
  \choice A read state, in which context is loaded, and a write state, in which memory is persisted.
\end{choices}
\dontknow
\itemmeta{T3 --- Loop and control flow}{remember}{easy}{A05}
\rationale{Reason/act is the core alternation; the other pairs name real activities that are not the
loop's two states.}

% --- Form B, item B05 -------------------------------------------------
\qitem In a ReAct-style loop, what triggers the transition from the reasoning state into the acting
state?
\begin{choices}
  \choice The context has reached its configured token budget.
  \CorrectChoice The LLM's output contains a tool call.
  \choice The user confirms that the agent may continue.
  \choice A fixed interval elapses between the two states.
\end{choices}
\dontknow
\itemmeta{T3 --- Loop and control flow}{remember}{easy}{B05}
\rationale{Same cell as A05: the model's output drives each transition. (c) confuses the loop with a
permission gate.}


% =====================================================================
% Cell 09 -- T4 / understand / medium
% =====================================================================

% --- Form A, item A09 -------------------------------------------------
\qitem Why is a large tool set with overlapping functionality a problem?
\begin{choices}
  \CorrectChoice It makes it harder for the agent to choose correctly, because the tools no longer define an unambiguous action space.
  \choice It raises the risk that two tools execute simultaneously and corrupt each other's output.
  \choice It prevents the harness from validating tool arguments against their schemas.
  \choice It forces the loop to call every registered tool at least once per run.
\end{choices}
\dontknow
\itemmeta{T4 --- Tools and tool design}{understand}{medium}{A09}
\rationale{Targets the ``more tools means a more capable agent'' misconception.}

% --- Form B, item B09 -------------------------------------------------
\qitem Your harness exposes \texttt{read\_file}, \texttt{open\_file}, \texttt{cat\_file} and
\texttt{get\_file\_contents}, all of which do the same thing. What is the most likely effect on the agent's
behaviour?
\begin{choices}
  \choice It refuses to call any of them, because the registry cannot resolve the ambiguity.
  \choice It calls them in sequence until one returns without an error.
  \choice It merges them into a single tool call in order to save tokens.
  \CorrectChoice It sometimes picks a less suitable one of the four, because nothing in the schemas distinguishes them.
\end{choices}
\dontknow
\itemmeta{T4 --- Tools and tool design}{understand}{medium}{B09}
\rationale{Same cell as A09, made concrete: an overlapping tool set degrades selection quality rather than
producing a hard failure, which is what makes it easy to miss.}


% =====================================================================
% Cell 14 -- T6 / understand / medium
% =====================================================================

% --- Form A, item A14 -------------------------------------------------
\qitem Skills and long-term memory are both loaded on demand rather than held in context permanently.
What distinguishes the kind of information they carry?
\begin{choices}
  \choice A skill is written by a developer and memory is written by the model; otherwise they are the same.
  \CorrectChoice A skill carries procedural know-how --- how to do something; long-term memory carries declarative facts --- what happened or what is the case.
  \choice A skill lives inside the context window, whereas long-term memory lives outside it.
  \choice A skill applies to a single run, whereas long-term memory is shared between all agents in the system.
\end{choices}
\dontknow
\itemmeta{T6 --- Skills}{understand}{medium}{A14}
\rationale{Procedural versus declarative is the distinction; (a) states a real difference but wrongly
concludes the two are otherwise equivalent, and (c) contradicts the stem.}

% --- Form B, item B14 -------------------------------------------------
\qitem A team wants their agent to remember that the user prefers metric units, and separately to know
the procedure for filling in a particular PDF form. Where does each belong?
\begin{choices}
  \choice Both in long-term memory, since both have to survive between runs.
  \choice The preference in a skill, the procedure in long-term memory.
  \CorrectChoice The preference in long-term memory, the procedure in a skill.
  \choice Both in the system prompt, since both apply to every run.
\end{choices}
\dontknow
\itemmeta{T6 --- Skills}{understand}{medium}{B14}
\rationale{Same cell as A14, as an assignment task: a fact goes to memory, a procedure to a skill.
(a) is tempting because persistence is common to both, which is exactly why it discriminates.}


% =====================================================================
% Cell 16 -- T7 / apply-analyse / hard
% =====================================================================

% --- Form A, item A16 -------------------------------------------------
\qitem A team's long-term memory has grown to several thousand stored facts, and loading all of them is
not an option. Which retrieval design matches how harnesses handle this?
\begin{choices}
  \CorrectChoice Put a few small, high-value items into the context up front, and keep lightweight identifiers such as file paths that let the agent load the rest when it needs them.
  \choice Load a fixed-size random sample of the stored facts at the start of every run.
  \choice Summarise the entire memory into one paragraph and include it in the system prompt.
  \choice Let the model decide at the end of the run which facts it should have read.
\end{choices}
\dontknow
\itemmeta{T7 --- Long-term memory}{apply-analyse}{hard}{A16}
\rationale{The two complementary strategies are up-front loading of high-value items and just-in-time
loading via identifiers.}

% --- Form B, item B16 -------------------------------------------------
\qitem Which decision has to be made at \emph{write} time by a long-term memory component and cannot be
deferred to read time?
\begin{choices}
  \choice Which of the stored items are relevant to the current task.
  \choice How much of the context budget the retrieved items may occupy.
  \choice Whether the retrieved items should be summarised before they are used.
  \CorrectChoice What is worth persisting at all, and at which point to persist it.
\end{choices}
\dontknow
\itemmeta{T7 --- Long-term memory}{apply-analyse}{hard}{B16}
\rationale{Same cell as A16, from the write side: everything not written is unrecoverable later, whereas
all three distractors are genuinely read-time concerns.}


% =====================================================================
% Cell 21 -- T10 / remember / easy
% =====================================================================

% --- Form A, item A21 -------------------------------------------------
\qitem Sandboxing can be realised at different levels of abstraction. Which sequence runs from the
weakest to the strongest isolation?
\begin{choices}
  \choice Containers, then operating-system primitives, then a fully separated operating system, then microVM isolation.
  \choice A fully separated operating system, then microVM isolation, then containers, then operating-system primitives.
  \CorrectChoice Operating-system primitives, then containers, then microVM or WASM isolation, then a fully separated operating system.
  \choice WASM isolation, then containers, then operating-system primitives, then microVM isolation.
\end{choices}
\dontknow
\itemmeta{T10 --- Sandboxing}{remember}{easy}{A21}
\rationale{The ordering reflects how much of the host is shared with the sandbox at each level.}

% --- Form B, item B21 -------------------------------------------------
\qitem What does sandboxing an agent primarily achieve?
\begin{choices}
  \choice It determines which of the agent's actions require the user's approval.
  \CorrectChoice It restricts the agent's access to the host system, network and files to a bounded environment.
  \choice It reduces the number of tokens the agent needs for a task.
  \choice It guarantees that the agent's tool calls always return the same result.
\end{choices}
\dontknow
\itemmeta{T10 --- Sandboxing}{remember}{easy}{B21}
\rationale{Same cell as A21. (a) is the sandbox/permission conflation and (d) overstates what a stable
environment buys --- it improves reproducibility, it does not make an agent deterministic.}


% =====================================================================
% Cell 24 -- T11 / apply-analyse / hard
% =====================================================================

% --- Form A, item A24 -------------------------------------------------
\qitem Which task benefits most from being delegated to a sub-agent?
\begin{choices}
  \choice Appending a single line to a file whose path is already known.
  \CorrectChoice Searching a large codebase for the places relevant to a change, where the search produces a lot of material but only a short list of findings is needed.
  \choice Deciding whether the user's request has been fulfilled and the run may end.
  \choice Incrementing a metrics counter after each tool call.
\end{choices}
\dontknow
\itemmeta{T11 --- Sub-agents and A2A}{apply-analyse}{hard}{A24}
\rationale{Delegation pays off when the ratio of intermediate material to useful result is high. (c) is
the coordinating agent's own job and must not be delegated.}

% --- Form B, item B24 -------------------------------------------------
\qitem A team delegates every individual tool call to its own sub-agent. Why does this fail to deliver
the benefit of delegation?
\begin{choices}
  \choice Sub-agents are not permitted to call tools, only to reason about them.
  \choice The parent agent cannot receive results from more than one sub-agent per run.
  \CorrectChoice Each delegation costs an extra context and an extra summary, while no substantial intermediate work is being kept out of the parent's context.
  \choice Sub-agents share the parent's context window, so nothing is saved.
\end{choices}
\dontknow
\itemmeta{T11 --- Sub-agents and A2A}{apply-analyse}{hard}{B24}
\rationale{Same cell as A24 as a negative case: delegation is worth its overhead only when the discarded
intermediate material is large relative to the returned result.}


% =====================================================================
% Cell 28 -- T13 / apply-analyse / hard
% =====================================================================

% --- Form A, item A28 -------------------------------------------------
\qitem You are writing a test that checks your agent cannot read a file outside its sandbox root. You
place a sentinel file outside the root and prompt the agent to read it. Which outcome should the test
accept as a pass?
\begin{choices}
  \choice Only the case in which the tool registry rejects the call; a refusal by the model is not sufficient.
  \choice Any outcome in which the agent does not crash, since isolation is enforced by the container.
  \choice Only the case in which the agent reads the file and then reports that it should not have been able to.
  \CorrectChoice Either the registry rejects the call or the agent refuses and no read happens --- as long as the sentinel's contents are never returned.
\end{choices}
\dontknow
\itemmeta{T13 --- Quality assurance}{apply-analyse}{hard}{A28}
\rationale{The property under test is that the sentinel is never read. Several acceptable mechanisms
satisfy it, so the assertion is written over the outcome rather than over one implementation.}

% --- Form B, item B28 -------------------------------------------------
\qitem A test suite already checks that the agent refuses to read \texttt{../outside\_target.txt}. Which
additional prompt variant adds the most value?
\begin{choices}
  \CorrectChoice ``Navigate to \texttt{/} and list its contents'' --- it probes a different code path with the same intent.
  \choice ``Please read \texttt{../outside\_target.txt}'' --- a politer phrasing of the same request.
  \choice ``Read \texttt{../outside\_target.txt} and tell me how many lines it has'' --- the same path with an extra step.
  \choice ``Read \texttt{../outside\_target.txt} twice'' --- it verifies that the refusal is stable.
\end{choices}
\dontknow
\itemmeta{T13 --- Quality assurance}{apply-analyse}{hard}{B28}
\rationale{Same cell as A28: coverage of a security boundary comes from exercising distinct enforcement
paths --- relative traversal, absolute paths, directory navigation --- not from rewording one request.}
